% !TEX program = xelatex
\documentclass[11pt]{article}
\usepackage{fontspec}
\usepackage{xeCJK}
\setmainfont{CMU Serif}
\newcommand{\setcjkserifandmono}[1]{%
  \setCJKmainfont{#1}%
  \setCJKmonofont{#1}%
}
\IfFontExistsTF{Noto Serif CJK SC}{\setcjkserifandmono{Noto Serif CJK SC}}{%
  \IfFontExistsTF{Source Han Serif SC}{\setcjkserifandmono{Source Han Serif SC}}{%
    \setcjkserifandmono{FandolSong-Regular}%
  }%
}%
\usepackage[margin=1in]{geometry}
\usepackage{amsmath}
\usepackage{booktabs}

\setlength{\parindent}{0pt}
\setlength{\parskip}{0.3\baselineskip}
\setlength{\emergencystretch}{2em}
\setlength{\textfloatsep}{5pt plus 1pt minus 1pt}
\setlength{\floatsep}{4pt plus 1pt minus 1pt}
\setlength{\intextsep}{4pt plus 1pt minus 1pt}
\setlength{\abovecaptionskip}{2pt}
\setlength{\belowcaptionskip}{1pt}
\setlength{\tabcolsep}{3pt}
\renewcommand{\arraystretch}{0.88}

\begin{document}

{\LARGE MiniOneRec SID 研究阶段性汇报}

\vspace{0.6\baselineskip}

\textbf{2026 年 5 月 6 日}

\section{Motivation}

原版 MiniOneRec 只使用了纯语义信息嵌入向量（item title + description）训练 RQ-VAE，构造三层 SID（$z^{(1)}, z^{(2)}, z^{(3)}$），每层 256 个码，下游 LLM 自回归预测目标 SID。

核心问题：\textbf{语义相近 $\neq$ 协同相近}。两个商品描述可能非常相似，但用户购买场景和使用人群完全不同。我们想做的事是——把图结构承载的协同信息，按层级感知的方式融入 SID 构建中，构造一个更容易被下游 LLM 学习的 SID 码本空间。

\section{Baseline}

原版 MiniOneRec 在  Industrial\_and\_Scientific 上的复现结果。

\begin{center}
\begin{tabular}{lrrrrrrrr}
\toprule
& \multicolumn{4}{c}{NDCG} & \multicolumn{4}{c}{HR} \\
\cmidrule(lr){2-5} \cmidrule(lr){6-9}
& @1 & @3 & @5 & @10 & @1 & @3 & @5 & @10 \\
\midrule
Baseline SFT & 0.06706 & 0.08501 & 0.09315 & \textbf{0.10372} & 0.06706 & 0.09839 & 0.11824 & 0.15089 \\
Baseline RL  & 0.07324 & 0.08903 & 0.09704 & \textbf{0.10726} & 0.07324 & 0.10038 & 0.11979 & 0.15133 \\
\bottomrule
\end{tabular}
\end{center}

\section{协同图信号注入}

不同的 SID 层级承担不同角色，应该接收不同粒度的协同图信号。

\medskip
\textbf{层级感知图正则}\\
引入三张图——$\mathbf{A}_{\rm coarse}$（序列共现粗图）、$\mathbf{A}_{\rm mid}$（谱分解中尺度图）、$\mathbf{A}_{\rm local}$（局部转移细图），分别挂到 RQ-VAE 的三层做图平滑：

\[
\mathcal{L}_{\rm v1} = \mathcal{L}_{\rm sem} + \lambda_c\mathcal{L}_{\rm graph}^{(1)} + \lambda_m\mathcal{L}_{\rm graph}^{(2)} + \lambda_l\mathcal{L}_{\rm graph}^{(3)},
\]

$\lambda_c=0.05$, $\lambda_m=0.15$, $\lambda_l=0.05$。L1 接粗图，L2 接中图，L3 接局部图。问题：固定全局权重对所有物品同等对待——本已稳定的语义结构被协同图过度修正。

\medskip
\textbf{歧义感知加权}\\
保留三层图分配，但引入了歧义先验 $a_i \in [0,1]$，用于衡量一个物品的语义结构和协同结构之间的不一致程度。歧义高的物品加强图监督权重，歧义低的物品加强语义保持（使其不被协同图过度改写）。同时引入语义近邻图作为语义结构锚点。

该方法 SFT NDCG@10 = 0.10271，RL NDCG@10 = 0.10432。没有超过 Baseline。

\section{图载体升级（TAGCF / FaGSP / MGDCF / Seq2Graph）}

动机：怀疑前述方法的谱中图不够好，尝试更复杂的图构造。

\begin{itemize}
  \item \textbf{TAGCF（属性增强中图）}：融合物品属性到拓扑图。NDCG@10 = 0.09758。
  \item \textbf{FaGSP（谱带滤波）}：用不同谱带代替均匀平滑。
  \item \textbf{MGDCF（多层图扩散）}：多层图代替单层图。
  \item \textbf{Seq2Graph-lite（高阶转移图）}：从序列历史构造高阶图。NDCG@10 = 0.09518。
\end{itemize}

结论：\textbf{换更复杂的图载体不能解决问题}。宽泛图信号容易破坏原版 SID 的可路由性。

\section{从图平滑到排序，以及选择性分离}

\medskip
\textbf{选择性推拉}\\
我们尝试在 L2 同时做拉近正样本和推开弱负样本（中层推拉），但 NDCG@10 = 0.09261。问题在于推开机制太粗——把所有非邻居当负样本，但实际上"图上没边"可能只是未观测到的交互。

受 CoST 启发，我们将 L2 目标改为对比 InfoNCE（正样本来自中图，负样本来自"语义近但协同弱"对），NDCG@10 = 0.09719。

\medskip
\textbf{Collab-ranking 排序对比系列}\\
不再用图平滑损失，而是把 L2 写成显式排序问题——$s_{ip}^{(2)} \ge s_{in}^{(2)} + m$，$p$ 是协同正样本，$n$ 是"语义近但协同弱"的困难负样本，$m=0.1$。

\begin{itemize}
  \item 局部多跳做正样本图：NDCG@10 = 0.09940。
  \item L1 额外用逆歧义加权做粗图拉近（与歧义感知相反：低歧义物品的语义结构更可靠，在粗路由层赋予更高权重）：NDCG@10 = 0.10094。
  \item K1=128（硬压缩 L1 容量）：NDCG@10 = 0.08955。
\end{itemize}

\section{最小编辑与精确推开}

回到最精简的思路：\textbf{只加最轻量的协同损失}。

\medskip
\textbf{最小协同编辑}\\
只在 L3 加局部协同拉近：NDCG@10 = 0.10159。\\
只在 L2 加局部多跳排序损失：NDCG@10 = 0.10165。

\medskip
\textbf{QCR 冲突推开}\\
只对\textbf{语义近 + 协同图弱连 + 当前共享 L2 前缀}的样本推开。

Tokenizer 侧：冲突 11/3686，活跃 L1 码 117，唯一 L2 前缀 2632。但下游 NDCG@10 = 0.09981。

\section{关键下游结果汇总}

\begin{center}
\begin{tabular}{lrrrrrrrr}
\toprule
& \multicolumn{4}{c}{NDCG} & \multicolumn{4}{c}{HR} \\
\cmidrule(lr){2-5} \cmidrule(lr){6-9}
& @1 & @3 & @5 & @10 & @1 & @3 & @5 & @10 \\
\midrule
Baseline RL & 0.07324 & 0.08903 & 0.09704 & \textbf{0.10726} & 0.07324 & 0.10038 & 0.11979 & \textbf{0.15133} \\
Baseline SFT & 0.06706 & 0.08501 & 0.09315 & \textbf{0.10372} & 0.06706 & 0.09839 & 0.11824 & 0.15089 \\
\midrule
\multicolumn{9}{c}{\textit{协同图注入}} \\
歧义感知加权 SFT & 0.07059 & 0.08451 & 0.09253 & 0.10271 & 0.07059 & 0.09508 & 0.11471 & 0.14626 \\
歧义感知加权 RL & 0.07434 & 0.09054 & 0.09630 & 0.10432 & 0.07434 & 0.10280 & 0.11692 & 0.14185 \\
\midrule
\multicolumn{9}{c}{\textit{最小编辑与精确推开}} \\
L2 局部多跳 & 0.06618 & 0.08283 & 0.09144 & 0.10165 & 0.06618 & 0.09486 & 0.11582 & 0.14736 \\
L3 协同拉近 & 0.06684 & 0.08445 & 0.09226 & 0.10159 & 0.06684 & 0.09773 & 0.11692 & 0.14604 \\
Collab-ranking 逆歧义加权 & 0.06398 & 0.08397 & 0.09185 & 0.10094 & 0.06398 & 0.09883 & 0.11802 & 0.14604 \\
QCR 冲突推开 & 0.06728 & 0.08579 & 0.09219 & 0.09981 & 0.06728 & 0.09949 & 0.11516 & 0.13876 \\
\midrule
\multicolumn{9}{c}{\textit{其他代表性变体}} \\
TAGCF 属性增强中图 & 0.07103 & 0.08505 & 0.08946 & 0.09758 & 0.07103 & 0.09530 & 0.10611 & 0.13148 \\
CoST 启发对比量化 & 0.06706 & 0.08149 & 0.08859 & 0.09719 & 0.06706 & 0.09221 & 0.10942 & 0.13611 \\
Seq2Graph 高阶转移图 & 0.06530 & 0.08132 & 0.08778 & 0.09518 & 0.06530 & 0.09354 & 0.10920 & 0.13236 \\
\bottomrule
\end{tabular}
\end{center}

\section{AttnRQ-Identity}

原版 RQ-VAE 训练目标是重建损失加量化损失：

\[
\mathcal{L}_{\rm rec} = \frac{1}{Nd}\sum_{i=1}^N \|\hat{x}_i - x_i\|_2^2,\qquad
\mathcal{L}_{\rm rq} = \frac{1}{3}\sum_{l=1}^3 \Big( \|\mathrm{sg}[r_i^{(l)}] - q_i^{(l)}\|_2^2 + 0.25\|r_i^{(l)} - \mathrm{sg}[q_i^{(l)}]\|_2^2 \Big),
\]
\[
\mathcal{L}_{\rm original} = \mathcal{L}_{\rm rec} + \mathcal{L}_{\rm rq}.
\]

其中 $x_i$ 是物品的语义嵌入，$r_i^{(l)}$ 是第 $l$ 层残差，$q_i^{(l)}$ 是其量化向量，$\mathrm{sg}[\cdot]$ 为 stop-gradient。重建路径使用三层残差码的固定求和：

\[
h_i = q_i^{(1)} + q_i^{(2)} + q_i^{(3)},\qquad \hat{x}_i = f_{\rm dec}(h_i).
\]

我们将固定求和替换为恒等初始化的注意力加权

\[
\tilde h_i = \sum_{l=1}^{3}\gamma_i^{(l)} q_i^{(l)},\qquad
\gamma_i^{(l)} = 3 \cdot {\rm softmax}_l(w_l^\top {\rm RMSNorm}(q_i^{(l)})),
\]
\[
\mathcal{L}_{\rm ours} = \|\hat{x}_i - x_i\|_2^2 + \mathcal{L}_{\rm rq},\qquad \hat{x}_i = f_{\rm dec}(\tilde h_i).
\]

$w_l$ 全零初始化，因此初始时 $\gamma_i^{(1)}=\gamma_i^{(2)}=\gamma_i^{(3)}=1$，模型严格等价于原版 RQ-VAE。训练过程中 $w_l$ 可学习，允许不同物品自适应调节各残差层的贡献度。

当前进度：

\begin{center}
\begin{tabular}{lrrrl}
\toprule
& 训练损失 & 原始冲突 & 生成冲突 & 平均 $\gamma$ \\
\midrule
原版 RQ-VAE 基线 & 0.20211 & 0.09929 & 16/3686 & -- \\
AttnRQ & 0.26489 & 0.12127 & 13/3686 & [1.010, 0.995, 0.995] \\
\bottomrule
\end{tabular}
\end{center}

方法训练稳定，$\gamma$ 接近 1，权重学到轻微偏置（略依赖 L1），生成冲突小幅改善。下游 SFT 尚未执行。

\end{document}
